\documentclass[12pt,a4paper]{article}

\usepackage[margin=2.5cm]{geometry}
\usepackage{times}
\usepackage{graphicx}
\usepackage{amsmath}
\usepackage{booktabs}
\usepackage{xcolor}
\usepackage{mdframed}
\usepackage{enumitem}
\usepackage{titlesec}
\usepackage{hyperref}
\usepackage{microtype}

\graphicspath{{figures/}}

% Colour scheme
\definecolor{speakblue}{RGB}{30, 90, 160}
\definecolor{notegray}{RGB}{100, 100, 100}
\definecolor{highlightyellow}{RGB}{255, 248, 200}
\definecolor{sectionred}{RGB}{180, 30, 30}

% Section formatting
\titleformat{\section}{\large\bfseries\color{sectionred}}{}{0em}{}[\titlerule]
\titleformat{\subsection}{\normalsize\bfseries\color{speakblue}}{}{0em}{}

% Speaker box
\newmdenv[
  backgroundcolor=highlightyellow,
  linecolor=speakblue,
  linewidth=1.5pt,
  leftmargin=0pt,
  rightmargin=0pt,
  innerleftmargin=10pt,
  innerrightmargin=10pt,
  innertopmargin=6pt,
  innerbottommargin=6pt,
]{speakbox}

% Note box
\newmdenv[
  backgroundcolor=gray!10,
  linecolor=notegray,
  linewidth=0.8pt,
  leftmargin=0pt,
  rightmargin=0pt,
  innerleftmargin=10pt,
  innerrightmargin=10pt,
  innertopmargin=5pt,
  innerbottommargin=5pt,
]{notebox}

\newcommand{\say}[1]{%
  \begin{speakbox}
  \textbf{\color{speakblue}[SPEAK]:} #1
  \end{speakbox}
  \vspace{4pt}
}

\newcommand{\note}[1]{%
  \begin{notebox}
  \textit{\color{notegray}[NOTE]: #1}
  \end{notebox}
  \vspace{4pt}
}

\newcommand{\action}[1]{%
  \vspace{2pt}
  \noindent\textbf{\color{sectionred}$\blacktriangleright$ ACTION:} \textit{#1}
  \vspace{4pt}
}

\begin{document}

\begin{center}
  {\LARGE \textbf{Demo Presentation Script}}\\[6pt]
  {\large Accent-Robust ASR using Wav2Vec 2.0 with Adversarial Training}\\[4pt]
  {\normalsize Karan Reddy (B22AI023) \quad Ale Anwesh (B22AI005) \quad Abhijan Theja (B22AI025)}\\[4pt]
  {\small Speech \& Audio Processing --- IIT Jodhpur}
\end{center}

\vspace{4pt}
\begin{notebox}
\textit{[NOTE]: Yellow boxes are what you say out loud. Gray boxes are private reminders.
Action lines tell you what to show/click on screen. Estimated total time: \textbf{12--15 minutes}.}
\end{notebox}

\vspace{8pt}

%=============================================================================
\section{Part 1 --- Introduction (2 minutes)}
%=============================================================================

\action{Open the GitHub repository page in the browser.}

\say{Good morning ma'am. We are team B22AI023, B22AI005, and B22AI025, and our project is on
\textbf{Accent-Robust Automatic Speech Recognition} using Wav2Vec 2.0 with Adversarial Training.}

\say{The core problem we are solving is this --- modern ASR systems like Google Speech or Alexa
work very well for native American English speakers, but they fail badly for speakers with
Indian, Arabic, or European accents. In some cases the Word Error Rate goes from 3\% for a
native speaker all the way up to 45\% for a non-native speaker. That is a serious fairness
problem.}

\say{Our solution combines two ideas. First, we use \textbf{Wav2Vec 2.0}, which is a
state-of-the-art self-supervised speech model from Meta AI that already understands speech
very well. Second, we add a \textbf{Gradient Reversal Layer}, or GRL, which forces the encoder
to forget what accent the speaker has while still learning how to transcribe speech correctly.
The result is a model that performs well across all accent groups, not just the majority.}

%=============================================================================
\section{Part 2 --- Architecture (2 minutes)}
%=============================================================================

\action{Open \texttt{report/report.tex} or the PDF and point to the architecture diagram.}

\say{Here is our system architecture. Raw audio at 16 kHz goes into the Wav2Vec 2.0 encoder.
This encoder has a CNN feature extractor that we keep frozen, and 12 Transformer layers that
we fine-tune.}

\say{The encoder produces hidden states. These go into two branches simultaneously.
The first branch is the \textbf{CTC Decoder} --- this is our ASR head that produces the
text transcription. The second branch is the \textbf{Accent Classifier}, which has two heads:
a supervised head that learns to identify accents, and an adversarial head that goes through
the Gradient Reversal Layer.}

\say{The GRL is mathematically very simple --- in the forward pass it is an identity function,
it does nothing. But during backpropagation it \textit{negates} the gradients. So when the
accent classifier tries to learn accent features, the negated gradients push the encoder in
the opposite direction, forcing it to remove accent information from its representations.}

\say{Our total training loss is:
$\mathcal{L} = \mathcal{L}_{\text{CTC}} + \lambda_s \mathcal{L}_{\text{accent}}^{s}
+ \lambda_a \mathcal{L}_{\text{accent}}^{a}$, where we use $\lambda_s = 0.1$ and
$\lambda_a = 1.0$.}

%=============================================================================
\section{Part 3 --- Code Walkthrough (3 minutes)}
%=============================================================================

\action{Open VS Code or file explorer. Navigate to \texttt{src/models/}.}

\subsection{Gradient Reversal Layer}

\action{Open \texttt{src/models/gradient\_reversal.py}.}

\say{This is our Gradient Reversal Layer implementation. We use PyTorch's custom autograd
function to override the backward pass. The forward method returns the input unchanged.
The backward method returns the gradient multiplied by negative alpha. We also implement
Ganin et al.'s annealing schedule where alpha starts at 0 and gradually increases to 1
during training, so the adversarial signal does not destabilize early training.}

\subsection{Main Model}

\action{Open \texttt{src/models/adversarial\_asr.py}.}

\say{Here is the main model. We load the full \texttt{Wav2Vec2ForCTC} pretrained model from
HuggingFace so that we get the pretrained CTC head weights. Then we attach our Accent
Classifier with the GRL. The forward method returns three things: CTC logits for transcription,
supervised accent logits, and adversarial accent logits.}

\subsection{Training Loss}

\action{Open \texttt{src/training/losses.py}.}

\say{This is our combined loss function. The CTC loss handles transcription. The supervised
accent loss only runs on samples where we know the accent label. The adversarial accent loss
also only runs on labeled samples, but because the GRL has already negated the gradients,
adding this loss pushes the encoder to be accent-invariant. We add all three together --- the
sign reversal is handled inside the GRL, not here.}

\subsection{Accent-Stratified Sampling}

\action{Open \texttt{src/data/sampling.py}.}

\say{One more important component --- our accent-stratified sampler. Since real accent datasets
are very unbalanced --- there are far more American English samples than Arabic or Korean ---
a naive training loop would just optimize for the majority accent. Our sampler assigns each
sample a weight inversely proportional to how common its accent is, so every accent gets
equal representation in each training batch.}

%=============================================================================
\section{Part 4 --- Results (3 minutes)}
%=============================================================================

\action{Open \texttt{results/evaluation/} folder. Show the PNG files.}

\subsection{WER Comparison}

\action{Open \texttt{wer\_comparison.png} and \texttt{per\_accent\_wer.png}.}

\say{These are our main results. We have three models: the Baseline which is just
\texttt{wav2vec2-base-960h} evaluated directly, the Accent-Adapted Phase 1 model after
fine-tuning on accented data, and our Adversarial model from Phase 2.}

\say{The baseline gives an overall WER of 27.96\% with a $\Delta$WERmax of 24.73\%. That
24.73\% gap means the best-performing accent group and the worst-performing one are nearly
25 percentage points apart --- very unfair. After our full adversarial training pipeline,
overall WER comes down to 18.15\% and $\Delta$WERmax drops to 16.10\%. That is a
\textbf{35\% relative WER reduction} and a \textbf{35\% fairness gap reduction}.}

\say{Looking at the per-accent chart, we can see the European accent had the highest baseline
WER at 41\%, and it comes down the most. American English barely changes, which confirms
that our model is not sacrificing performance on native speakers to help non-native ones ---
it is genuinely improving for everyone.}

\subsection{Training Curves}

\action{Open \texttt{training\_curves.png}.}

\say{These are the training curves across all three phases. You can see Phase 0 converges
quickly on the baseline. Phase 1 accent adaptation continues to reduce WER. Phase 2 with
the GRL shows further reduction in both overall WER and in the fairness gap.}

%=============================================================================
\section{Part 5 --- Responsible AI (2 minutes)}
%=============================================================================

\action{Open \texttt{report/figures/intersectional\_fairness.png}.}

\say{We also did a comprehensive Responsible AI analysis. This intersectional fairness
heatmap breaks down WER by both accent and gender. The left panel is the Baseline --- you
can see female speakers in non-native accent groups like European and Arabic have the
highest WER, because they face compounded bias from both their accent and their gender.}

\say{The center panel shows our adversarial model, where WER is uniformly lower across
all groups. The right panel shows the absolute improvement --- the greener the cell, the
more we improved. Female non-native speakers benefit the most, which is exactly what we
want a fairness-aware model to do.}

\say{We also implemented Grad-CAM explainability over the Wav2Vec 2.0 encoder layers.
In the Baseline model, the accent classifier's Grad-CAM activations are very strong at
consonant cluster positions --- for example, retroflex consonants in Indian English or
vowel elongation in Arabic English. After adversarial training, these activations are
significantly reduced and more spread out, which confirms that the GRL is successfully
removing accent-specific cues from the encoder's upper layers.}

%=============================================================================
\section{Part 6 --- Closing (1 minute)}
%=============================================================================

\action{Open the GitHub repository page again.}

\say{To summarize: we built an end-to-end accent-robust ASR system using Wav2Vec 2.0 and
a Gradient Reversal Layer. We trained in three phases --- baseline, accent adaptation, and
adversarial training. Our model achieves a 35\% reduction in WER and a 35\% reduction in
the $\Delta$WERmax fairness gap. We also performed intersectional analysis and Grad-CAM
explainability to ensure our results are transparent and interpretable.}

\say{All our code, configs, results, and the CVPR-format report are available on GitHub.
Everything is reproducible with a single script using seed 42. Thank you ma'am, we are
happy to take any questions.}

\vspace{8pt}
\begin{notebox}
\textbf{Likely questions and answers:}
\begin{itemize}[leftmargin=1.5em]
  \item \textit{``Why did you freeze the CNN feature extractor?''} --- The CNN layers
  extract low-level acoustic features like spectral patterns. These are universal across
  accents, so fine-tuning them would risk destroying the pretrained representations.
  Only the Transformer layers need to be adapted.

  \item \textit{``Why not just collect more accented data?''} --- Data collection is
  expensive and slow, especially for low-resource accent groups. Our approach uses
  domain-adversarial training to get accent invariance from the existing data, without
  needing additional collection.

  \item \textit{``How does the GRL not hurt ASR performance?''} --- The CTC loss and
  the adversarial loss optimize for different representations simultaneously. The CTC
  loss ensures phonetic information is preserved while the GRL removes only the
  accent-specific signal. The supervised accent head ($\lambda_s = 0.1$) provides
  a soft regularization that prevents the GRL from going too far and destroying
  useful acoustic features.

  \item \textit{``What is $\Delta$WERmax?''} --- It is the difference between the
  highest per-accent WER and the lowest per-accent WER. A lower value means the
  model is more fair across accent groups.

  \item \textit{``Why Wav2Vec 2.0 and not Whisper?''} --- Wav2Vec 2.0 is a
  feature encoder, making it easier to attach the GRL at the representation level.
  Whisper is an encoder-decoder model, which would require a different adversarial
  training setup. Wav2Vec 2.0 is also better studied for domain adaptation tasks.
\end{itemize}
\end{notebox}

\end{document}
